% Fate's Edge SRD −−− Health, Fatigue, and Harm (Section 09)

\section{Health, Fatigue, and Harm}

\subsection{Fatigue Track}

Each character has a Fatigue track equal to their \textbf{Body} attribute (minimum 1). Fatigue represents exertion, strain, and the accumulation of minor physical toll.

\textbf{Marking Fatigue:} Mark Fatigue for the following:
\begin{itemize}
  ℹtem Performing a \textbf{Charge} or \textbf{Skirmish} action (1 Fatigue).
  ℹtem Forced marches, extreme effort, or environmental strain (GM discretion).
  ℹtem Magical costs (e.g., Pushing a Rite, Backlash, Obligation overflow).
  ℹtem Converting a Boon effect to Fatigue (substitution rule).
\end{itemize}

\subsection{Effect of Fatigue}

Each point of Fatigue worsens your \textbf{Position} by one step:
\[
\text{Dominant} → \text{Controlled} → \text{Desperate}
\]

If you are already at \textbf{Desperate Position}, each additional point of Fatigue instead imposes a \textbf{–1 die penalty} on all your rolls (cumulative).

\subsection{Harm Levels}

Harm represents serious injury. There are three levels:

\begin{tabular}{@{}ll@{}}
⊤rule
\textbf{Harm Level} & \textbf{Effect} \\
∣rule
Harm 1 & –1 die on actions directly related to the injury (GM discretion). \\
Harm 2 & –1 die on most actions until treated. \\
Harm 3 & Incapacitated or dying (cannot act; requires stabilization). \\
⊥tomrule
\end{tabular}

\subsection{Taking Harm and Fatigue −−− The Procedure}

When you take Harm from an attack, hazard, or magical backlash, resolve in the following order:

\begin{enumerate}
  ℹtem \textbf{Convert.} Use the \textit{Armor Conversion Table} to turn incoming Harm into Fatigue. Minimum 1 Fatigue per hit, even if conversion would give 0.
  ℹtem \textbf{Apply remaining Harm.} If any Harm remains after conversion:
    \begin{itemize}
      ℹtem Clear \textbf{all existing Fatigue} (the shock of a wound resets exhaustion).
      ℹtem Increase your Harm level by the number of Harm points remaining.
    \end{itemize}
  ℹtem \textbf{Add Fatigue.} Add the Fatigue from step 1 to your Fatigue track.
  ℹtem \textbf{Check for Overflow.} If your Fatigue total equals or exceeds your Body (track size):
    \begin{itemize}
      ℹtem Increase Harm by one level.
      ℹtem Clear all Fatigue.
      ℹtem Any excess Fatigue beyond the track size is lost.
    \end{itemize}
    Overflow can occur multiple times from a single hit (e.g., with Body 1 or 2).
\end{enumerate}

\subsection{Armor Conversion Table}

Armor does not reduce incoming Harm directly. Instead, it converts Harm to Fatigue according to the table below.

\begin{tabular}{@{}lccc@{}}
⊤rule
\textbf{Armor Type} & \textbf{Harm 1} & \textbf{Harm 2} & \textbf{Harm 3} \\
∣rule
\textbf{Light}   & Fatigue 1            & Harm 1 + Fatigue 1   & Harm 2 + Fatigue 1 \\
\textbf{Medium}  & Fatigue 1            & Fatigue 1            & Harm 2 + Fatigue 1 \\
\textbf{Heavy}   & Fatigue 1            & Fatigue 1            & Fatigue 2 \\
⊥tomrule
\end{tabular}

\textbf{Example:} Light armor vs Harm 2 becomes \textit{Harm 1 + Fatigue 1}. Heavy armor vs Harm 3 becomes \textit{Fatigue 2}.

\subsection{Recovery}

⫕section{Fatigue Recovery}
\begin{itemize}
  ℹtem \textbf{Short rest} (1 hour): Remove 2 Fatigue.
  ℹtem \textbf{Full night's rest} (8 hours): Remove all Fatigue.
\end{itemize}

⫕section{Harm Recovery}
\begin{itemize}
  ℹtem \textbf{Medical care:} A successful Medicine roll (DV 3–5, depending on severity) can reduce Harm by one level after a scene of treatment. A critical success may remove two levels.
  ℹtem \textbf{Magical healing:} A \texttt{[HEAL]} spell or Rite can reduce Harm by one level (see Magic system).
  ℹtem \textbf{Natural recovery:} Without medical or magical aid, one level of Harm heals after one week of rest. Harm 3 requires intervention.
\end{itemize}

\subsection{Substituting Fatigue for Boons}

If you have no Boons remaining (or wish to conserve them), you may pay 1 Fatigue in place of 1 Boon to activate any effect that normally costs a Boon. The GM may veto this substitution if the fiction makes it implausible (e.g., activating a delicate Asset while exhausted). This does not change the effect's outcome; it merely changes the resource spent.

\subsection{Example: Armor and Overflow}

\textit{Heavy armor, Body 2, currently Fatigue 1, Harm 0. Hit with Harm 3.}

\begin{enumerate}
  ℹtem Convert: Harm 3 $→$ Fatigue 2 (from table).
  ℹtem Apply remaining Harm: 0 remaining (all converted). No Harm applied, no Fatigue cleared.
  ℹtem Add Fatigue: 1 (current) + 2 = 3. Track size is 2 $→$ overflow.
  ℹtem Overflow: Harm 0 $→$ 1. Clear all Fatigue. Excess lost.
  ℹtem Result: Harm 1, Fatigue 0.
\end{enumerate}

\subsection{Death and Dying}

At \textbf{Harm 3}, a character is incapacitated. They cannot take actions. They may still speak weakly or crawl at the GM's discretion.

⫕section{Stabilization}
\begin{itemize}
  ℹtem An ally may use an action to stabilize the dying character with a successful Medicine roll (DV 3). On success, the character remains at Harm 3 but is stable (no further deterioration).
  ℹtem Without stabilization, the character will die at the end of the scene unless circumstances change (e.g., magical healing, arrival of a skilled healer).
\end{itemize}

⫕section{Raising Harm Above 3}
If a character at Harm 3 takes additional Harm, they die immediately. Resurrection or miraculous recovery is possible only through major rituals, Patron intervention, or narrative circumstances (GM discretion).

